% 2項真理関数 f_0〜f_15 のCNF（和積標準形）一覧
% 第7章「2項論理関数」の添字記法 f_{b1b2b3b4}（b1b2b3b4はV3,V2,V1,V0への出力）
% に基づき，p,qの2変項で作れる16個の真理関数すべてを最も単純なCNFで表す。
% 対策プリント問4ではf2,f6,f11,f13のみを扱ったので，本資料でf0〜f15を網羅する。
% 各f_iについて，真理値表→最大項→簡単化，という導出過程を省略なく書く。
\documentclass[autodetect-engine,dvi=dvipdfmx,ja=standard]{bxjsarticle}

\makeatletter
\providecommand*{\input@path}{}
\g@addto@macro\input@path{{./}{../includes/}}
\makeatother

\input{protocol.tex}
\input{preamble.tex}
\usepackage{longtable}

\title{2項真理関数 \(f_0\)〜\(f_{15}\) のCNF一覧（導出過程つき）}
\author{論理学 第7章 対応}
\date{\today}

\begin{document}
\maketitle

第7章の記法にしたがい，\(p,q\) の2変項から作れる \(2^4=16\) 通りの真理関数
\(f_0,f_1,\ldots,f_{15}\) を，添字 \(f_{b_1b_2b_3b_4}\)（入力
\((V_3,V_2,V_1,V_0)=\bigl((1,1),(1,0),(0,1),(0,0)\bigr)\) に対する出力を順に並べたビット列）
で表す。本資料は，対策プリント問4で扱った \(f_2,f_6,f_{11},f_{13}\) を含む16個すべてについて，
\textbf{真理値表→最大項→簡単化}という導出過程を省略せずに書き，最も単純な和積標準形（CNF）を求める。

\section{全体表（参照用）}

\begin{center}
\scriptsize
\setlength{\tabcolsep}{3pt}
\begin{tabular}{c|cc|cccccccccccccccc}
\toprule
 & \(p\) & \(q\)
 & \(f_0\) & \(f_1\) & \(f_2\) & \(f_3\) & \(f_4\) & \(f_5\) & \(f_6\) & \(f_7\)
 & \(f_8\) & \(f_9\) & \(f_{10}\) & \(f_{11}\) & \(f_{12}\) & \(f_{13}\) & \(f_{14}\) & \(f_{15}\)
\\
\midrule
\(V_3\) & 1 & 1 & 0&0&0&0&0&0&0&1&1&1&1&1&1&1&1&1 \\
\(V_2\) & 1 & 0 & 0&0&0&0&1&1&1&1&0&0&0&1&1&1&1&1 \\
\(V_1\) & 0 & 1 & 0&0&1&1&0&0&1&1&0&0&1&0&0&1&1&1 \\
\(V_0\) & 0 & 0 & 0&1&0&1&0&1&0&1&0&1&0&1&0&1&0&1 \\
\bottomrule
\end{tabular}
\end{center}

\section{導出に使う道具}

CNFは，真理値表で\textbf{偽になる行}から作った最大項の連言として求める
（\(x\)の値が偽の行では\(x\)自身，真の行では\(\sim x\)をリテラルに採る——対策プリント問3の
「別解」で使ったのと同じ方法）。\(p,q\)の2変項しかないので，最大項は多くても4個で，
以下の2つの恒等式だけで必ず1〜2節に単純化できる。

\[
  \text{恒等式①（併合）：}\quad(x\lor A)\land(\sim x\lor A)\equiv A
\]
\[
  \text{恒等式②（吸収）：}\quad A\land(B\lor\sim A)\equiv A\land B
\]
恒等式①は分配律から，恒等式②も分配律から
（\(A\land(B\lor\sim A)\equiv(A\land B)\lor(A\land\sim A)\equiv(A\land B)\lor\bot\equiv A\land B\)）
それぞれ直接出る。

\section{各関数の導出}

\subsection*{\(f_0=f_{0000}\)（矛盾・\(\bot\)）}
すべての行で偽。4つの最大項\((\sim p\lor\sim q),(\sim p\lor q),(p\lor\sim q),(p\lor q)\)ができる。
最初の2つは恒等式①（共通部分\(\sim p\)）で\(\sim p\)に，残り2つも恒等式①（共通部分\(p\)）で\(p\)にまとまり，
\(\sim p\land p\)——変項とその否定が同時に真になることはないので恒偽式である。
\textbf{CNF：}\(p\land\sim p\)

\subsection*{\(f_1=f_{0001}\)（NOR：\(\sim(p\lor q)\)）}
真になるのは\((p,q)=(0,0)\)のみ。偽の3行から最大項
\((\sim p\lor\sim q),(\sim p\lor q),(p\lor\sim q)\)ができる。
最初の2つは共通部分\(\sim p\)で恒等式①により\(\sim p\)にまとまる。
残る\((p\lor\sim q)\)との連言\(\sim p\land(p\lor\sim q)\)に恒等式②（\(A=\sim p,B=\sim q\)）を適用すると
\(\sim p\land\sim q\)。
\textbf{CNF：}\(\sim p\land\sim q\)

\subsection*{\(f_2=f_{0010}\)（\(\sim p\land q\)）}
真になるのは\((p,q)=(0,1)\)のみ。偽の3行から最大項
\((\sim p\lor\sim q),(\sim p\lor q),(p\lor q)\)ができる。
最初の2つは共通部分\(\sim p\)で\(\sim p\)にまとまる。
残る\((p\lor q)\)との連言\(\sim p\land(p\lor q)\)に恒等式②（\(A=\sim p,B=q\)）を適用すると\(\sim p\land q\)。
\textbf{CNF：}\(\sim p\land q\)

\subsection*{\(f_3=f_{0011}\)（射影：\(\sim p\)）}
真になるのは\((p,q)=(0,1),(0,0)\)。偽の2行から最大項\((\sim p\lor\sim q),(\sim p\lor q)\)ができ，
共通部分\(\sim p\)により恒等式①でそのまま\(\sim p\)にまとまる。
\textbf{CNF：}\(\sim p\)

\subsection*{\(f_4=f_{0100}\)（非含意：\(p\land\sim q\)）}
真になるのは\((p,q)=(1,0)\)のみ。偽の3行から最大項
\((\sim p\lor\sim q),(p\lor\sim q),(p\lor q)\)ができる。
後ろ2つは共通部分\(p\)で\(p\)にまとまる。
残る\((\sim p\lor\sim q)\)との連言\(p\land(\sim p\lor\sim q)\)に恒等式②（\(A=p,B=\sim q\)）を適用すると
\(p\land\sim q\)。
\textbf{CNF：}\(p\land\sim q\)

\subsection*{\(f_5=f_{0101}\)（射影：\(\sim q\)）}
真になるのは\((p,q)=(1,0),(0,0)\)。偽の2行から最大項\((\sim p\lor\sim q),(p\lor\sim q)\)ができ，
共通部分\(\sim q\)により恒等式①でそのまま\(\sim q\)にまとまる。
\textbf{CNF：}\(\sim q\)

\subsection*{\(f_6=f_{0110}\)（XOR：\(p\oplus q\)）}
真になるのは\((p,q)=(1,0),(0,1)\)。偽の2行は\((1,1)\)と\((0,0)\)で，最大項は
\((\sim p\lor\sim q)\)と\((p\lor q)\)。この2つは\(p,q\)どちらのリテラルも極性が逆
（\(p\)と\(\sim q\)、\(\sim p\)と\(q\)がそれぞれ入れ替わっている）で，恒等式①が要求する
「1変項だけ違って他は共通」という形になっていない。したがってこれ以上まとめられず，
2節構成のままが最も単純な形である。
\textbf{CNF：}\((\sim p\lor\sim q)\land(p\lor q)\)

\subsection*{\(f_7=f_{0111}\)（NAND：\(\sim(p\land q)\)）}
偽になるのは\((p,q)=(1,1)\)のみ。最大項は1個だけで，それがそのままCNFになる。
\textbf{CNF：}\(\sim p\lor\sim q\)

\subsection*{\(f_8=f_{1000}\)（AND：\(p\land q\)）}
真になるのは\((p,q)=(1,1)\)のみ。偽の3行から最大項
\((\sim p\lor q),(p\lor\sim q),(p\lor q)\)ができる。
後ろ2つは共通部分\(p\)で\(p\)にまとまる。
残る\((\sim p\lor q)\)との連言\(p\land(\sim p\lor q)\)に恒等式②（\(A=p,B=q\)）を適用すると\(p\land q\)。
\textbf{CNF：}\(p\land q\)

\subsection*{\(f_9=f_{1001}\)（XNOR：\(p\equiv q\)）}
真になるのは\((p,q)=(1,1),(0,0)\)。偽の2行は\((1,0)\)と\((0,1)\)で，最大項は
\((\sim p\lor q)\)と\((p\lor\sim q)\)。\(f_6\)と同様，2つとも\(p,q\)両方のリテラルが入れ替わっており，
恒等式①の形にならないので併合できない。
\textbf{CNF：}\((\sim p\lor q)\land(p\lor\sim q)\)

\subsection*{\(f_{10}=f_{1010}\)（射影：\(q\)）}
真になるのは\((p,q)=(1,1),(0,1)\)。偽の2行から最大項\((\sim p\lor q),(p\lor q)\)ができ，
共通部分\(q\)により恒等式①でそのまま\(q\)にまとまる。
\textbf{CNF：}\(q\)

\subsection*{\(f_{11}=f_{1011}\)（含意：\(p\supset q\)）}
偽になるのは\((p,q)=(1,0)\)のみ。最大項は1個だけで，それがそのままCNFになる。
\textbf{CNF：}\(\sim p\lor q\)

\subsection*{\(f_{12}=f_{1100}\)（射影：\(p\)）}
真になるのは\((p,q)=(1,1),(1,0)\)。偽の2行から最大項\((p\lor\sim q),(p\lor q)\)ができ，
共通部分\(p\)により恒等式①でそのまま\(p\)にまとまる。
\textbf{CNF：}\(p\)

\subsection*{\(f_{13}=f_{1101}\)（逆向きの含意：\(q\supset p\)）}
偽になるのは\((p,q)=(0,1)\)のみ。最大項は1個だけで，それがそのままCNFになる。
\textbf{CNF：}\(p\lor\sim q\)

\subsection*{\(f_{14}=f_{1110}\)（OR：\(p\lor q\)）}
偽になるのは\((p,q)=(0,0)\)のみ。最大項は1個だけで，それがそのままCNFになる。
\textbf{CNF：}\(p\lor q\)

\subsection*{\(f_{15}=f_{1111}\)（トートロジー・\(\top\)）}
すべての行で真であり，偽になる行が0個なので作られる最大項も0個。
最大項の連言（CNF）は空，すなわち恒真式である。
\textbf{CNF：}\(p\lor\sim p\)

\section{結果一覧（まとめ）}

\begin{center}
\begin{longtable}{c|c|l|l}
\toprule
\(f_i\) & \(f_{b_1b_2b_3b_4}\) & 最も単純なCNF & 名称・備考 \\
\midrule
\endhead
\(f_0\) & \(f_{0000}\) & \(p\land\sim p\)（恒偽） & 矛盾（\(\bot\)） \\
\(f_1\) & \(f_{0001}\) & \(\sim p\land\sim q\) & NOR：\(\sim(p\lor q)\) \\
\(f_2\) & \(f_{0010}\) & \(\sim p\land q\) & 非含意の逆 \\
\(f_3\) & \(f_{0011}\) & \(\sim p\) & 射影：\(\sim p\) \\
\(f_4\) & \(f_{0100}\) & \(p\land\sim q\) & 非含意：\(p\land\sim q\) \\
\(f_5\) & \(f_{0101}\) & \(\sim q\) & 射影：\(\sim q\) \\
\(f_6\) & \(f_{0110}\) & \((\sim p\lor\sim q)\land(p\lor q)\) & XOR（排他的選言）：\(p\oplus q\) \\
\(f_7\) & \(f_{0111}\) & \(\sim p\lor\sim q\) & NAND：\(\sim(p\land q)\) \\
\(f_8\) & \(f_{1000}\) & \(p\land q\) & AND（連言）：\(p\land q\) \\
\(f_9\) & \(f_{1001}\) & \((\sim p\lor q)\land(p\lor\sim q)\) & XNOR：\(p\equiv q\) \\
\(f_{10}\) & \(f_{1010}\) & \(q\) & 射影：\(q\) \\
\(f_{11}\) & \(f_{1011}\) & \(\sim p\lor q\) & 含意：\(p\supset q\) \\
\(f_{12}\) & \(f_{1100}\) & \(p\) & 射影：\(p\) \\
\(f_{13}\) & \(f_{1101}\) & \(p\lor\sim q\) & 逆向きの含意：\(q\supset p\) \\
\(f_{14}\) & \(f_{1110}\) & \(p\lor q\) & OR（選言）：\(p\lor q\) \\
\(f_{15}\) & \(f_{1111}\) & \(p\lor\sim p\)（恒真） & トートロジー（\(\top\)） \\
\bottomrule
\end{longtable}
\end{center}

補足として，\(f_i\)と\(f_{15-i}\)は常に互いに否定の関係にある（各行の出力ビットがすべて反転するため）。
\(f_0/f_{15}\)（矛盾/トートロジー），\(f_1/f_{14}\)（NOR/OR），\(f_2/f_{13}\)，\(f_3/f_{12}\)，
\(f_4/f_{11}\)，\(f_5/f_{10}\)，\(f_6/f_9\)（XOR/XNOR），\(f_7/f_8\)（NAND/AND）が対になっている。

\end{document}
